\documentclass[10pt,a4paper]{amsart}
\usepackage{amsmath,amssymb,amsthm,graphicx}
\usepackage{verbatim}
\newenvironment{reference}{\comment}{\endcomment}
%\newenvironment{reference}{}{}
\newenvironment{slogan}{\comment}{\endcomment}
\newenvironment{history}{\comment}{\endcomment}

% For commutative diagrams we use Xy-pic
\usepackage[all]{xy}

% We use 2cell for 2-commutative diagrams.
\xyoption{2cell}
% Theorem environments.
%
\theoremstyle{plain}
\newtheorem{theorem}[subsection]{Theorem}
\newtheorem{proposition}[subsection]{Proposition}
\newtheorem{lemma}[subsection]{Lemma}

\theoremstyle{definition}
\newtheorem{definition}[subsection]{Definition}
\newtheorem{example}[subsection]{Example}
\newtheorem{exercise}[subsection]{Exercise}
\newtheorem{situation}[subsection]{Situation}

\theoremstyle{remark}
\newtheorem{remark}[subsection]{Remark}
\newtheorem{remarks}[subsection]{Remarks}

\numberwithin{equation}{subsection}

% Macros
%
\def\lim{\mathop{\mathrm{lim}}\nolimits}
\def\colim{\mathop{\mathrm{colim}}\nolimits}
\def\Spec{\mathop{\mathrm{Spec}}}
\def\Hom{\mathop{\mathrm{Hom}}\nolimits}
\def\Ext{\mathop{\mathrm{Ext}}\nolimits}
\def\SheafHom{\mathop{\mathcal{H}\!\mathit{om}}\nolimits}
\def\SheafExt{\mathop{\mathcal{E}\!\mathit{xt}}\nolimits}
\def\Sch{\mathit{Sch}}
\def\Mor{\mathop{\mathrm{Mor}}\nolimits}
\def\Ob{\mathop{\mathrm{Ob}}\nolimits}
\def\Sh{\mathop{\mathit{Sh}}\nolimits}
\def\NL{\mathop{N\!L}\nolimits}
\def\CH{\mathop{\mathrm{CH}}\nolimits}
\def\proetale{{pro\text{-}\acute{e}tale}}
\def\etale{{\acute{e}tale}}
\def\QCoh{\mathit{QCoh}}
\def\Ker{\mathop{\mathrm{Ker}}}
\def\Im{\mathop{\mathrm{Im}}}
\def\Coker{\mathop{\mathrm{Coker}}}
\def\Coim{\mathop{\mathrm{Coim}}}

% Boxtimes
%
\DeclareMathSymbol{\boxtimes}{\mathbin}{AMSa}{"02}

%
% Macros for moduli stacks/spaces
%
\def\QCohstack{\mathcal{QC}\!\mathit{oh}}
\def\Cohstack{\mathcal{C}\!\mathit{oh}}
\def\Spacesstack{\mathcal{S}\!\mathit{paces}}
\def\Quotfunctor{\mathrm{Quot}}
\def\Hilbfunctor{\mathrm{Hilb}}
\def\Curvesstack{\mathcal{C}\!\mathit{urves}}
\def\Polarizedstack{\mathcal{P}\!\mathit{olarized}}
\def\Complexesstack{\mathcal{C}\!\mathit{omplexes}}
% \Pic is the operator that assigns to X its picard group, usage \Pic(X)
% \Picardstack_{X/B} denotes the Picard stack of X over B
% \Picardfunctor_{X/B} denotes the Picard functor of X over B
\def\Pic{\mathop{\mathrm{Pic}}\nolimits}
\def\Picardstack{\mathcal{P}\!\mathit{ic}}
\def\Picardfunctor{\mathrm{Pic}}
\def\Deformationcategory{\mathcal{D}\!\mathit{ef}}
\usepackage{cleveref}
\crefname{originallemma}{original lemma}{original lemmas}
\Crefname{originallemma}{Original lemma}{Original lemmas}
\makeatletter
\newlabel{lemma-dimension-at-most-one}{{02P7}{source}{}{sourcepointer-1}{}}
\newlabel{lemma-dimension-at-most-one@cref}{{[originallemma][0][]02P7}{[1][0][]source}{}{}{}}
\makeatother
\crefname{originallemma}{original lemma}{original lemmas}
\Crefname{originallemma}{Original lemma}{Original lemmas}
\makeatletter
\newlabel{lemma-compare-det}{{02P8}{source}{}{sourcepointer-2}{}}
\newlabel{lemma-compare-det@cref}{{[originallemma][0][]02P8}{[1][0][]source}{}{}{}}
\makeatother
\title{For a finite length module over a local ring the determinant is one-dimensional and generated by the symbol of any spanning admissible sequence}
\author{The Stacks Project Authors}
\date{stacks-project:chow.tex:local-snapshot-20261012; GNU FDL 1.2 or later}
\begin{document}
\maketitle
\noindent\textit{Editorial note.} All delivered mathematical text is the author's own, copied from the pinned Stacks Project chapter file chow.tex; no statement or proof step was rewritten, reordered or removed.\par\medskip
\noindent\textit{Editorial note.} Preamble provenance: the pinned chapter file opens with an input of the shared preamble and carries no preamble text of its own, so the project's single Stacks-wide preamble staged at sources/stacks-chow/preamble.tex is excerpted verbatim as two spans, lines 14-18 (Xy-pic setup) and lines 154-217 (the native theorem declarations and mathematical macros). Nothing was added, removed or redefined; the wrapper class and the amsmath/amssymb/amsthm/graphicx base packages are supplied by the assembler.\par\medskip
\noindent\textit{Editorial note.} The counted statement and proof contain no cross reference, no citation and no included file, so the delivered body is self contained.\par\medskip
\noindent\textit{Editorial note.} Canonical location of the counted proposition: Stacks tag 02P9 (sources/stacks-chow/pinned-original-source.tex line 15850).\par\medskip
\noindent\textit{Allowed standard prerequisites (editorial).}\begin{itemize}
\item The determinant det\_kappa(M) of a finite length module over a local ring, its symbols and its admissible relations. Same chapter, the construction referred to as 'defined above' in the statement of this lemma. Fixes the object to be shown nonzero and the three families of relations that psi must respect.
\item det\_kappa(M) has dimension at most one and is generated by the symbol of a spanning admissible sequence. Same chapter, Lemma labelled lemma-dimension-at-most-one, cited at proof lines 15859-15862 together with its proof. Supplies the dimension bound and the generation statement, reducing the lemma to nonvanishing.
\item The length one case of the determinant. Same chapter, Lemma labelled lemma-compare-det, cited at proof line 15864. Base case of the induction on the length.
\end{itemize}\medskip
\ifdefined\hypertarget\hypertarget{sourcepointer-1}{}\fi
\noindent\textit{Source reference.} Original lemma 02P7: sources/stacks-chow/pinned-original-source.tex, source lines 15761--15761. Same-chapter pointer into the author environment lemma at source line 15761; invoked at proof line 15862 of the counted proof. It is retained only as a navigation pointer; every load-bearing step is inside the delivered span.\par\medskip
\ifdefined\hypertarget\hypertarget{sourcepointer-2}{}\fi
\noindent\textit{Source reference.} Original lemma 02P8: sources/stacks-chow/pinned-original-source.tex, source lines 15827--15827. Same-chapter pointer into the author environment lemma at source line 15827; invoked at proof line 15864 of the counted proof. It is retained only as a navigation pointer; every load-bearing step is inside the delivered span.\par\medskip
\begin{lemma}
\label{lemma-determinant-dimension-one}
Let $R$ be a local ring with maximal ideal $\mathfrak m$ and
residue field $\kappa$. Let $M$ be a finite length $R$-module.
The determinant $\det_\kappa(M)$ defined above is a $\kappa$-vector
space of dimension $1$. It is generated by the symbol
$[f_1, \ldots, f_l]$ for any admissible sequence such
that $\langle f_1, \ldots f_l \rangle = M$.
\end{lemma}

\begin{proof}
We know $\det_\kappa(M)$ has dimension at most $1$, and in fact that it
is generated by $[f_1, \ldots, f_l]$, by
Lemma \ref{lemma-dimension-at-most-one} and its proof.
We will show by induction on $l = \text{length}(M)$
that it is nonzero. For $l = 1$ it follows from Lemma \ref{lemma-compare-det}.
Choose a nonzero element $f \in M$
with $\mathfrak m f = 0$. Set $\overline{M} = M /\langle f \rangle$,
and denote the quotient map $x \mapsto \overline{x}$.
We will define a surjective map
$$
\psi : \det\nolimits_k(M) \to \det\nolimits_\kappa(\overline{M})
$$
which will prove the lemma since by induction the determinant of
$\overline{M}$ is nonzero.

\medskip\noindent
We define $\psi$ on symbols as follows.
Let $(e_1, \ldots, e_l)$ be an admissible sequence.
If $f \not \in \langle e_1, \ldots, e_l \rangle$ then
we simply set $\psi([e_1, \ldots, e_l]) = 0$.
If $f \in \langle e_1, \ldots, e_l \rangle$ then we choose
an $i$ minimal such that $f \in \langle e_1, \ldots, e_i \rangle$.
We may write $e_i = \lambda f + x$ for some unit $\lambda \in R$
and $x \in \langle e_1, \ldots, e_{i - 1} \rangle$.
In this case we set
$$
\psi([e_1, \ldots, e_l]) =
(-1)^i
\overline{\lambda}[\overline{e}_1, \ldots,
\overline{e}_{i - 1},
\overline{e}_{i + 1}, \ldots, \overline{e}_l].
$$
Note that it is indeed the case that
$(\overline{e}_1, \ldots,
\overline{e}_{i - 1},
\overline{e}_{i + 1}, \ldots, \overline{e}_l)$
is an admissible sequence in $\overline{M}$, so this makes sense.
Let us show that extending this rule $\kappa$-linearly to
linear combinations of symbols does indeed lead to a map on
determinants. To do this we have to show that the admissible
relations are mapped to zero.

\medskip\noindent
Type (a) relations. Suppose we have $(e_1, \ldots, e_l)$ an
admissible sequence and for some $1 \leq a \leq l$ we have
$e_a \in \langle e_1, \ldots, e_{a - 1}\rangle$.
Suppose that $f \in \langle e_1, \ldots, e_i\rangle$ with $i$ minimal.
Then $i \not = a$ and
$\overline{e}_a \in \langle \overline{e}_1, \ldots,
\hat{\overline{e}_i}, \ldots, \overline{e}_{a - 1}\rangle$ if $i < a$
or
$\overline{e}_a \in \langle \overline{e}_1, \ldots,
\overline{e}_{a - 1}\rangle$ if $i > a$.
Thus the same admissible relation for $\det_\kappa(\overline{M})$ forces
the symbol $[\overline{e}_1, \ldots,
\overline{e}_{i - 1},
\overline{e}_{i + 1}, \ldots, \overline{e}_l]$
to be zero as desired.

\medskip\noindent
Type (b) relations. Suppose we have $(e_1, \ldots, e_l)$ an
admissible sequence and for some $1 \leq a \leq l$ we have
$e_a = \lambda e'_a + x$ with $\lambda \in R^*$, and
$x \in \langle e_1, \ldots, e_{a - 1}\rangle$.
Suppose that $f \in \langle e_1, \ldots, e_i\rangle$ with $i$ minimal.
Say $e_i = \mu f + y$ with $y \in \langle e_1, \ldots, e_{i - 1}\rangle$.
If $i < a$ then the desired equality is
$$
(-1)^i
\overline{\lambda}
[\overline{e}_1,
\ldots,
\overline{e}_{i - 1},
\overline{e}_{i + 1},
\ldots,
\overline{e}_l]
=
(-1)^i
\overline{\lambda}
[\overline{e}_1,
\ldots,
\overline{e}_{i - 1},
\overline{e}_{i + 1},
\ldots,
\overline{e}_{a - 1},
\overline{e}'_a,
\overline{e}_{a + 1},
\ldots,
\overline{e}_l]
$$
which follows from $\overline{e}_a = \lambda \overline{e}'_a + \overline{x}$
and the corresponding admissible relation for $\det_\kappa(\overline{M})$.
If $i > a$ then the desired equality is
$$
(-1)^i
\overline{\lambda}
[\overline{e}_1,
\ldots,
\overline{e}_{i - 1},
\overline{e}_{i + 1},
\ldots,
\overline{e}_l]
=
(-1)^i
\overline{\lambda}
[\overline{e}_1,
\ldots,
\overline{e}_{a - 1},
\overline{e}'_a,
\overline{e}_{a + 1},
\ldots,
\overline{e}_{i - 1},
\overline{e}_{i + 1},
\ldots,
\overline{e}_l]
$$
which follows from $\overline{e}_a = \lambda \overline{e}'_a + \overline{x}$
and the corresponding admissible relation for $\det_\kappa(\overline{M})$.
The interesting case is when $i = a$. In this case we have
$e_a = \lambda e'_a + x = \mu f + y$. Hence also
$e'_a = \lambda^{-1}(\mu f + y - x)$. Thus we see that
$$
\psi([e_1, \ldots, e_l])
= (-1)^i \overline{\mu}
[\overline{e}_1, \ldots,
\overline{e}_{i - 1},
\overline{e}_{i + 1}, \ldots, \overline{e}_l]
=
\psi(
\overline{\lambda}
[e_1, \ldots, e_{a - 1}, e'_a, e_{a + 1}, \ldots, e_l]
)
$$
as desired.

\medskip\noindent
Type (c) relations. Suppose that $(e_1, \ldots, e_l)$
is an admissible sequence and
$\mathfrak m e_a \subset \langle e_1, \ldots, e_{a - 2}\rangle$.
Suppose that $f \in \langle e_1, \ldots, e_i\rangle$ with $i$ minimal.
Say $e_i = \lambda f + x$ with $x \in \langle e_1, \ldots, e_{i - 1}\rangle$.
We distinguish $4$ cases:

\medskip\noindent
Case 1: $i < a - 1$. The desired equality is
\begin{align*}
& (-1)^i
\overline{\lambda}
[\overline{e}_1,
\ldots,
\overline{e}_{i - 1},
\overline{e}_{i + 1},
\ldots,
\overline{e}_l] \\
& =
(-1)^{i + 1}
\overline{\lambda}
[\overline{e}_1,
\ldots,
\overline{e}_{i - 1},
\overline{e}_{i + 1},
\ldots,
\overline{e}_{a - 2},
\overline{e}_a,
\overline{e}_{a - 1},
\overline{e}_{a + 1},
\ldots,
\overline{e}_l]
\end{align*}
which follows from the type (c) admissible relation for
$\det_\kappa(\overline{M})$.

\medskip\noindent
Case 2: $i > a$. The desired equality is
\begin{align*}
& (-1)^i
\overline{\lambda}
[\overline{e}_1,
\ldots,
\overline{e}_{i - 1},
\overline{e}_{i + 1},
\ldots,
\overline{e}_l] \\
& =
(-1)^{i + 1}
\overline{\lambda}
[\overline{e}_1,
\ldots,
\overline{e}_{a - 2},
\overline{e}_a,
\overline{e}_{a - 1},
\overline{e}_{a + 1},
\ldots,
\overline{e}_{i - 1},
\overline{e}_{i + 1},
\ldots,
\overline{e}_l]
\end{align*}
which follows from the type (c) admissible relation for
$\det_\kappa(\overline{M})$.

\medskip\noindent
Case 3: $i = a$. We write $e_a = \lambda f + \mu e_{a - 1} + y$
with $y \in \langle e_1, \ldots, e_{a - 2}\rangle$. Then
$$
\psi([e_1, \ldots, e_l]) =
(-1)^a
\overline{\lambda}
[\overline{e}_1,
\ldots,
\overline{e}_{a - 1},
\overline{e}_{a + 1},
\ldots,
\overline{e}_l]
$$
by definition. If $\overline{\mu}$ is nonzero, then we have
$e_{a - 1} = - \mu^{-1} \lambda f + \mu^{-1}e_a - \mu^{-1} y$
and we obtain
$$
\psi(-[e_1, \ldots, e_{a - 2}, e_a, e_{a - 1}, e_{a + 1}, \ldots, e_l]) =
(-1)^a
\overline{\mu^{-1}\lambda}
[\overline{e}_1,
\ldots,
\overline{e}_{a - 2},
\overline{e}_a,
\overline{e}_{a + 1},
\ldots,
\overline{e}_l]
$$
by definition. Since in $\overline{M}$ we have
$\overline{e}_a = \mu \overline{e}_{a - 1} + \overline{y}$ we see
the two outcomes are equal by relation (a) for $\det_\kappa(\overline{M})$.
If on the other hand $\overline{\mu}$ is zero, then we can write
$e_a = \lambda f + y$ with $y \in \langle e_1, \ldots, e_{a - 2}\rangle$
and we have
$$
\psi(-[e_1, \ldots, e_{a - 2}, e_a, e_{a - 1}, e_{a + 1}, \ldots, e_l]) =
(-1)^a
\overline{\lambda}
[\overline{e}_1,
\ldots,
\overline{e}_{a - 1},
\overline{e}_{a + 1},
\ldots,
\overline{e}_l]
$$
which is equal to $\psi([e_1, \ldots, e_l])$.

\medskip\noindent
Case 4: $i = a - 1$. Here we have
$$
\psi([e_1, \ldots, e_l]) =
(-1)^{a - 1}
\overline{\lambda}
[\overline{e}_1,
\ldots,
\overline{e}_{a - 2},
\overline{e}_a,
\ldots,
\overline{e}_l]
$$
by definition. If $f \not \in \langle e_1, \ldots, e_{a - 2}, e_a \rangle$
then
$$
\psi(-[e_1, \ldots, e_{a - 2}, e_a, e_{a - 1}, e_{a + 1}, \ldots, e_l]) =
(-1)^{a + 1}\overline{\lambda}
[\overline{e}_1,
\ldots,
\overline{e}_{a - 2},
\overline{e}_a,
\ldots,
\overline{e}_l]
$$
Since $(-1)^{a - 1} = (-1)^{a + 1}$ the two expressions are the same.
Finally, assume $f \in \langle e_1, \ldots, e_{a - 2}, e_a \rangle$.
In this case we see that $e_{a - 1} = \lambda f + x$ with
$x \in \langle e_1, \ldots, e_{a - 2}\rangle$ and
$e_a = \mu f + y$ with $y \in \langle e_1, \ldots, e_{a - 2}\rangle$
for units $\lambda, \mu \in R$.
We conclude that both
$e_a \in \langle e_1, \ldots, e_{a - 1} \rangle$ and
$e_{a - 1} \in \langle e_1, \ldots, e_{a - 2}, e_a\rangle$.
In this case a relation of type (a) applies to both
$[e_1, \ldots, e_l]$ and
$[e_1, \ldots, e_{a - 2}, e_a, e_{a - 1}, e_{a + 1}, \ldots, e_l]$
and the compatibility of $\psi$ with these shown above to see that both
$$
\psi([e_1, \ldots, e_l])
\quad\text{and}\quad
\psi([e_1, \ldots, e_{a - 2}, e_a, e_{a - 1}, e_{a + 1}, \ldots, e_l])
$$
are zero, as desired.

\medskip\noindent
At this point we have shown that $\psi$ is well defined, and all that remains
is to show that it is surjective. To see this let
$(\overline{f}_2, \ldots, \overline{f}_l)$ be an admissible sequence
in $\overline{M}$. We can choose lifts $f_2, \ldots, f_l \in M$, and
then $(f, f_2, \ldots, f_l)$ is an admissible sequence in $M$.
Since $\psi([f, f_2, \ldots, f_l]) = [f_2, \ldots, f_l]$ we win.
\end{proof}

\end{document}
